\documentclass[twocolumn,aps,prxlife,longbibliography,amsmath,amssymb,floatfix,superscriptaddress,showkeys]{revtex4-2}
%  see todo_list.txt, more TODOs inline
\usepackage[colorlinks=true, linkcolor=blue, citecolor=blue, urlcolor=blue]{hyperref}
\usepackage{graphicx}

% \usepackage[monochrome]{xcolor} % ALL COLORS TURNED BLACK
\usepackage{xcolor} % WITH COLORS
\usepackage{bm}
\definecolor{lightgray}{rgb}{0.8,0.8,0.8}

\newcommand{\red}[1]{\textcolor{red}{#1}}

\begin{document}
\title{Rotation}

\date{}
\author{Maximilian Kotz}
\affiliation{Cluster of Excellence Physics of Life, TU Dresden, Dresden, Germany}
% \author{Veikko F. Geyer}
% \affiliation{B CUBE, TU Dresden, Dresden, Germany}


% \email{maximilian.kotz@tu-dresden.de}
\date{\today}

\begin{abstract}
\dots
\end{abstract}

\maketitle

\section{Methods used by others}
The persistence length $L_p$ (phenomenological quantity) is defined by the
\begin{equation}
  \langle \bm{t}(0)\cdot\bm{t}(s) \rangle = \langle \cos\left[ \theta(0) - \theta(s) \right] \rangle = C(s) = e^{-|s|/2L_p} \, .
  \label{eq:def_lp}
\end{equation} 

\subsection{Used in Graham et al.~\cite{graham_multi-platform_2014}}
\paragraph{Cosine correlation.}
One can directly determine $L_p$ by fitting an exponential to numerically estimated values of $C(s)$. Because of noise the values will be underestimated.
The authors~\cite{graham_multi-platform_2014} suggest with reference to~\cite{isambert_flexibility_1995,gittes_flexural_1993} to fit an exponential with free amplitude.
\red{fitted in log or real, fitted up to which s, fitted with which weight}

\paragraph{End to end distance.} 
The end-to-end distance as function of contour length is given by~\cite{howard2001mechanics,landau1980statistical}
\begin{equation}
  \langle R^2 \rangle = 2L_p^2 \left( e^{-L/L_p} - 1 + L/L_p \right)
\end{equation}
They used only actual end-to-end distance and therefore one datapoint per filament.
\red{cant we calculate all R2 depending on separation}

\paragraph{Fourier analysis.} 
From equipartition theorem we know \cite{gittes_flexural_1993}
\begin{equation}
  L_p = \frac{L^2}{n^2 \pi^2 \operatorname{Var}(a_n)}
\end{equation}

\subsection{Used in Wisanpitayakorn et al.~\cite{wisanpitayakorn_measurement_2022}}
\paragraph{Curvature.}
The setup is a chain with points sampled approximately equidistant with distance $\Delta s$.
For small $\Delta s$ one assumes the probability of a curvature at a given node ($\kappa=\theta/\Delta s$)
\begin{equation}
  P(\kappa) = \sqrt{\frac{L_p \Delta s}{2 \pi}} e^{-L_p \Delta s \kappa^2/2} \, .
\end{equation}
By taking just every mth node one gets chains with less points, there porbability distribution is given by
\begin{equation}
  P(\kappa^{(m)})=\sqrt{\frac{\Lambda^{(m)}}{2\pi}}e^{-\Lambda^{(m)} (\kappa^{(m)})^2/2}
\end{equation}
with \begin{equation}
  \Lambda^{(m)} = L_p \mu^{(m)} \Delta s^{(m)}
  \label{eq:lambda_Lp}
\end{equation}
and
\begin{equation}
  \mu^{(m)} = \frac{m}{(2m - 1)(1 - 1/m)/3 + 1}
\end{equation}
One estimates for each m the $\Lambda$ by fitting the cumulative distribution
\begin{equation}
  C(\kappa)^{(m)} = \frac{1}{2}\left[1 +  \operatorname{erf}\left( \sqrt{\frac{\Lambda}{2}}(\kappa^{(m)} - \kappa_0^{(m)})\right)\right]
\end{equation}
where $\kappa_0$ is the mean if the distribution.
Finally $L_p$ is obtained by the linear fit using Eq. \ref{eq:lambda_Lp}.

\paragraph{Fourier analysis.}
They used
\begin{equation}
  \langle (\frac{a_n - a_n^0}{L})^2 \rangle = \frac{1}{\pi^2 L_p n^2} + \langle \frac{4 \epsilon^2}{L^3}[1 + (N-1) \sin^2(\frac{n\pi}{2N})] \rangle
\end{equation}
 Where $\epsilon_k$ is random displacement and $N$ number of sampling points. 

\subsection{Used in van den Heuvel et al.~\cite{van_den_heuvel_persistence_2007}}
Gliding assay. The free end of the microtubule can bind to kinesin. It binds with a distance $d_i$ from the last binding. 
It dies so under an angle standart distributed with $\operatorname{Var}(\theta_i) = d_i/L_p$. Therefore after N connections we have
$\operatorname{Var}(\theta)= \sum_i \operatorname{Var}(\theta_i)$ and a total distance of $d=\sum_i d_i$, which gives
\begin{equation}
  \operatorname{Var}(\theta) = d/L_p
\end{equation}
They determined the variance for different lag times and performed linear fit. Ones has to be careful as the data is correlated.
Finally they performed linear regression.
They obtained $\theta$ from finite difference of consecutive frame of the tip of the MT.

\subsection{Used in Lamour et al.~\cite{lamour_easyworm_2014}}
\paragraph{Persistence length calculation}
Same as Cosine correlation in ~\cite{graham_multi-platform_2014}

\paragraph{End to end distance.}
Also same as in ~\cite{graham_multi-platform_2014} but here maybe in fact with taking also parts of the chain and not only the total one.

\paragraph{Secant}
They cite ~\cite{gittes_flexural_1993}.
Only valid for $L_p$ much longer than contour length.
Take a look at there page 2 fig 2.

\bibliography{rotation}


\end{document}
